\documentclass{article}
\usepackage{amsmath, amssymb, amsthm}
\usepackage{geometry}
\usepackage[UTF8]{ctex} % 用于中文支持，如果不需要中文，可以移除此行

\newtheorem{problem}{问题}
\newtheorem{solution}{解答}
\title{数值分析第三周作业}
\date{\today}
\author{李佩优PB23000270}
\geometry{a4paper,scale=0.8}
\begin{document}
\maketitle
\section*{6.3}

\subsection*{1}
用推广的牛顿均差方法求一个满足下列表值的二次多项式：

\begin{table}[h!]
  \begin{center}
    \begin{tabular}{c|c|c|c} % <-- Alignments: 1st column left, 2nd middle and 3rd right, with vertical lines in between
      x & 0 & 1 & 2\\
      \hline
        $p(x)$ & 2 & -4 & 44 \\
        $p'(x)$ & -9 & 4 & - \\
    \end{tabular}
  \end{center}
\end{table}

用推广的牛顿均差方法，构造重节点差商表。节点序列为 \(0,0,1,1,2\)，对应函数值 \(2,2,-4,-4,44\)，导数值 \(-9\) 和 \(4\) 用于重节点差商。计算各阶差商：

一阶：\(f[0,0]=-9\)，\(f[0,1]=-6\)，\(f[1,1]=4\)，\(f[1,2]=48\)

二阶：\(f[0,0,1]=3\)，\(f[0,1,1]=10\)，\(f[1,1,2]=44\)

三阶：\(f[0,0,1,1]=7\)，\(f[0,1,1,2]=17\)

四阶：\(f[0,0,1,1,2]=5\)

得到牛顿插值多项式：
\[
p(x)=2-9x+3x^2+7x^2(x-1)+5x^2(x-1)^2=5x^4-3x^3+x^2-9x+2.
\]
经验证满足所有表值。

\subsection*{2}
求一个满足上题中表值，并且满足$p(3)=2$的五次多项式$p$,

增加条件 \(p(3)=2\)，在节点序列后添加 \(3\)，对应函数值 \(2\)。继续计算差商： 

一阶：\(f[2,3]=-42\)

二阶：\(f[1,2,3]=-45\)

三阶：\(f[1,1,2,3]=-44.5\)

四阶：\(f[0,1,1,2,3]=-20.5\)

五阶：\(f[0,0,1,1,2,3]=-8.5\)

得五次多项式：
\[
p(x)=5x^4-3x^3+x^2-9x+2-8.5x^2(x-1)^2(x-2).
\]
展开为
\[
p(x)=-\frac{17}{2}x^5+39x^4-\frac{91}{2}x^3+18x^2-9x+2.
\]

\subsection*{3}
试求满足下列表值的次数最低的多项式$p$的公式：
$p(x_i)=y_i,\quad p'(x_i)=0 (0\leqslant i\leqslant n)$

设 \(l_i(x)=\prod_{j\neq i}\frac{x-x_j}{x_i-x_j}\)，则满足 \(p(x_i)=y_i\) 且 \(p'(x_i)=0\) 的次数最低的多项式为
\[
p(x)=\sum_{i=0}^n y_i\left[1-2l_i'(x_i)(x-x_i)\right]l_i^2(x),
\]
其中 \(l_i'(x_i)=\sum_{j\neq i}\frac{1}{x_i-x_j}\)。该多项式次数不超过 \(2n+1\)。

\subsection*{4}
如果插值问题 
$p(x_i)=c_{i0} p''(x_i) = c_{i2} (i=0,1)$
有一个三次多项式的解(对任意的$c_{ij}$)，试问对节点$x_0$和$x_1$应设置什么条件？


给定条件 \(p(x_i)=c_{i0}\)，\(p''(x_i)=c_{i2}\)（\(i=0,1\)），设三次多项式 \(p(x)=ax^3+bx^2+cx+d\)，代入得线性方程组。系数矩阵行列式为 \(12(x_0-x_1)^2\)。对任意常数 \(c_{ij}\) 有解当且仅当行列式非零，即 \(x_0\neq x_1\)。因此节点必须互异。

\(p(x)=5x^4-3x^3+x^2-9x+2\) \quad

\(p(x)=-\dfrac{17}{2}x^5+39x^4-\dfrac{91}{2}x^3+18x^2-9x+2\) \quad

\(p(x)=\sum_{i=0}^n y_i\left[1-2l_i'(x_i)(x-x_i)\right]l_i^2(x)\) \quad

\(x_0\neq x_1\)

\section*{6.4}

\subsection*{5}
确定下列函数是否为一个二次样条函数：

$$ f(x)=\left\{
\begin{aligned}
&x &\quad & x\in (-\infty, 1)\\
&-\frac{1}{2}(2-x)^2 + \frac{3}{2}& \quad & x\in [1,2] \\
&\frac{3}{2}&\quad & x\in [2,\infty)
\end{aligned}
\right.
$$

原函数在每个区间上都是二次多项式。  

在节点 \(x=1\) 处：

左极限 \(f(1^-)=1\)，右值 \(f(1)=-\frac{1}{2}(1)^2+\frac{3}{2}=1\)，连续。

左导数 \(f'(1^-)=1\)，右导数 \(f'(1)=2-1=1\)，一阶导数连续。

在节点 \(x=2\) 处：

左极限 \(f(2^-)=-\frac{1}{2}(0)^2+\frac{3}{2}=\frac{3}{2}\)，右值 \(f(2)=\frac{3}{2}\)，连续。

左导数 \(f'(2^-)=2-2=0\)，右导数 \(f'(2)=0\)，一阶导数连续。

因此，函数满足二次样条的条件（连续且一阶导数连续），故它是一个二次样条函数。

\subsection*{6}
上题中的函数是否是一个三次样条函数:

对于同一个函数，检查二阶导数：
第一段：\(f''(x)=0\)（\(x<1\)）
第二段：\(f''(x)=-1\)（\(1<x<2\)）
第三段：\(f''(x)=0\)（\(x>2\)）
在 \(x=1\) 处，左二阶导数为 \(0\)，右二阶导数为 \(-1\)，不连续；在 \(x=2\) 处，左二阶导数为 \(-1\)，右二阶导数为 \(0\)，也不连续。因此不满足三次样条要求的二阶导数连续，故不是三次样条函数。


\subsection*{7}
确定$a,b,c,d,e$的值，使得下列函数是一个三次样条：
$$ f(x)=\left\{
\begin{aligned}
&a(x-2)^2 +b(x-1)^3&\quad & x\in (-\infty, 1]\\
&c(x-2)^2& \quad & x\in [1,3] \\
&d(x-2)^2 +e(x-3)^3&\quad & x\in [3,\infty)
\end{aligned}
\right.
$$  
其次，确定这些参数的值，使得这个三次样条插值下列表值：
$f(0)=26,\,\,f(1)=7,\,\,f(4)=25\,\,$

(1)
三次样条，即在节点 \(x=1\) 和 \(x=3\) 处满足函数值、一阶导数、二阶导数连续。
计算各段导数：
- \(f_1(x)=a(x-2)^2+b(x-1)^3\)，\(f_1'(x)=2a(x-2)+3b(x-1)^2\)，\(f_1''(x)=2a+6b(x-1)\)
- \(f_2(x)=c(x-2)^2\)，\(f_2'(x)=2c(x-2)\)，\(f_2''(x)=2c\)
- \(f_3(x)=d(x-2)^2+e(x-3)^3\)，\(f_3'(x)=2d(x-2)+3e(x-3)^2\)，\(f_3''(x)=2d+6e(x-3)\)
在 \(x=1\) 处：
\[
f_1(1)=a,\quad f_2(1)=c \Rightarrow a=c
\]
\[
f_1'(1)=-2a,\quad f_2'(1)=-2c \Rightarrow a=c
\]
\[
f_1''(1)=2a,\quad f_2''(1)=2c \Rightarrow a=c
\]
在 \(x=3\) 处：
\[
f_2(3)=c,\quad f_3(3)=d \Rightarrow c=d
\]
\[
f_2'(3)=2c,\quad f_3'(3)=2d \Rightarrow c=d
\]
\[
f_2''(3)=2c,\quad f_3''(3)=2d \Rightarrow c=d
\]
因此，三次样条条件等价于 \(a=c=d\)，而 \(b\) 和 \(e\) 可任意。

(2)
再利用插值条件：
\(f(0)=26\)：\(f_1(0)=a(0-2)^2+b(0-1)^3=4a-b=26\)
\(f(1)=7\)：\(f_1(1)=a=7\)
\(f(4)=25\)：\(f_3(4)=d(4-2)^2+e(4-3)^3=4d+e=25\)
代入 \(a=7\)，得 \(4\times7-b=26\Rightarrow b=2\)；由 \(a=c=d=7\)，得 \(4\times7+e=25\Rightarrow e=-3\)。
故所求参数为：
\[
a=7,\quad b=2,\quad c=7,\quad d=7,\quad e=-3.
\]

\subsection*{11}
确定$a,b,c$的值，使得下列函数是一个具有节点$0,1,2$の三次样条：

\[
f(x)=\begin{cases}
3+x-9x^2, & x\in[0,1] \\[2pt]
a+b(x-1)+c(x-1)^2+d(x-1)^3, & x\in[1,2]
\end{cases}
\]
其次，确定$d$使得$\int_{0}^{2}[f''(x)]^2  \,dx $达到极小。
最后，求$d$的一个值，使得$f''(2)=0$，并解释为什么这个值不同于前面所确定的值。

(1)

左端在 \(x = 1\) 处：\(f(1) = 3 + 1 - 9 = -5\)，右端 \(f(1) = a\)，得 \(a = -5\)。

一阶导数：左端 \(f'(x) = 1 - 18x\)，\(f'(1) = -17\)；右端 \(f'(x) = b + 2c(x-1) + 3d(x-1)^2\)，\(f'(1) = b\)，得 \(b = -17\)。

二阶导数：左端 \(f''(x) = -18\)，\(f''(1) = -18\)；右端 \(f''(x) = 2c + 6d(x-1)\)，\(f''(1) = 2c\)，得 \(2c = -18\)，即 \(c = -9\)。

因此，\(a = -5\)，\(b = -17\)，\(c = -9\)。

(2)

此时 \(f''(x)\) 在 \([0,1]\) 上为常数 \(-18\)，在 \([1,2]\) 上为 \(f''(x) = -18 + 6d(x-1)\)。积分：
\[
I(d) = \int_0^1 324 \, dx + \int_1^2 \left[-18 + 6d(x-1)\right]^2 dx = 324 + \int_0^1 (324 - 216d t + 36d^2 t^2) \, dt,
\]
其中 \(t = x-1\)。计算得：
\[
I(d) = 324 + \left[324t - 108d t^2 + 12d^2 t^3\right]_0^1 = 648 - 108d + 12d^2.
\]
这是关于 \(d\) 的二次函数，极小值点满足 \(I'(d) = -108 + 24d = 0\)，解得 \(d = 4.5\)。

(3)

\(f''(2) = -18 + 6d = 0\) 得 \(d = 3\)。该值不同于极小化积分得到的 \(d = 4.5\)，因为前者对应自然样条条件（端点二阶导数为零），而后者是在给定第一段函数形式下极小化曲率平方积分，二者属于不同的优化目标。

\subsection*{12}
确定下列函数是否为一个三次样条
\[
f(x)=\begin{cases}
x^3+x, & x\leqslant 0 \\[2pt]
x^3-x, & x\geqslant 0 
\end{cases}
\]

证明$\lim_{x\uparrow 0}f''(x)=\lim_{x\downarrow 0}f''(x)$

在 \(x = 0\) 处：

函数值连续：\(f(0) = 0\)。

一阶导数：左导数为 \(3x^2 + 1\)，\(f'(0^-) = 1\)；右导数为 \(3x^2 - 1\)，\(f'(0^+) = -1\)，不相等。

二阶导数：左导数为 \(6x\)，\(f''(0^-) = 0\)；右导数为 \(6x\)，\(f''(0^+) = 0\)，相等。

由于一阶导数不连续，该函数不是三次样条。但二阶导数在 \(x = 0\) 处连续，即 \(\lim_{x \uparrow 0} f''(x) = \lim_{x \downarrow 0} f''(x) = 0\)。

\subsection*{14}
确定下列函数是否为一个自然三次样条
\[
f(x)=\begin{cases}
2(x+1)+(x+1)^3, & x\in[-1,0] \\[2pt]
3+5x+3x^2, & x\in[0,1] \\[2pt]
11+11(x-1)+3(x-1)^2-(x-1)^3, & x\in[1,2]
\end{cases}
\]

函数：
\[
f(x) = \begin{cases}
2(x+1) + (x+1)^3, & x \in [-1, 0] \\
3 + 5x + 3x^2, & x \in [0, 1] \\
11 + 11(x-1) + 3(x-1)^2 - (x-1)^3, & x \in [1, 2]
\end{cases}
\]
检查节点 \(x = 0\) 和 \(x = 1\) 处的连续性及导数连续性：

(1)在 \(x = 0\) 处：

左：\(f(0) = 2 \cdot 1 + 1 = 3\)，右：\(f(0) = 3\)，连续。

一阶导数：左 \(f'(x) = 2 + 3(x+1)^2\)，\(f'(0) = 5\)；右 \(f'(x) = 5 + 6x\)，\(f'(0) = 5\)，连续。

二阶导数：左 \(f''(x) = 6(x+1)\)，\(f''(0) = 6\)；右 \(f''(x) = 6\)，\(f''(0) = 6\)，连续。

(2)在 \(x = 1\) 处：

左：\(f(1) = 3 + 5 + 3 = 11\)，右：\(f(1) = 11\)，连续。

一阶导数：左 \(f'(x) = 5 + 6x\)，\(f'(1) = 11\)；右 \(f'(x) = 11 + 6(x-1) - 3(x-1)^2\)，\(f'(1) = 11\)，连续。

二阶导数：左 \(f''(x) = 6\)，\(f''(1) = 6\)；右 \(f''(x) = 6 - 6(x-1)\)，\(f''(1) = 6\)，连续。

(3)端点条件：

在 \(x = -1\) 处，\(f''(-1) = 6(-1+1) = 0\)。

在 \(x = 2\) 处，\(f''(2) = 6 - 6(2-1) = 0\)。

函数满足所有条件，是一个自然三次样条。

\subsection*{19}
试求出一个自然三次样条函数，它的结点是$-1,0,1$
并且取值是$f(-1)=5,\,\,f(0)=7,\,\,f(1)=9$

由于三点共线（斜率均为 \(2\)），自然三次样条（端点二阶导数为零）即为线性函数，故  
\[
S(x)=2x+7.
\]

\subsection*{26}
确定$(a,b,c)$的值，使得函数
\[
f(x)=\begin{cases}
x^3, & x\in[0,1] \\[2pt]
\frac{1}{2}(x-1)^3 + a(x-1)^2+b(x-1)+c, & x\in[1,3]
\end{cases}
\]

要使其成为三次样条（即 \(C^2\) 连续），需在 \(x=1\) 处满足函数值、一阶导、二阶导连续。

由 \(f(1^-)=1\)，\(f(1^+)=c\)，得 \(c=1\)。

由 \(f'(1^-)=3\)，\(f'(1^+)=b\)，得 \(b=3\)。

由 \(f''(1^-)=6\)，\(f''(1^+)=2a\)，得 \(a=3\)。

因此，\((a,b,c)=(3,3,1)\,\,\,S(x)=2x+7\)。

\end{document}
% \[
% f(x)=\begin{cases}
% a(x-2)^2+b(x-1)^3, & x\in(-\infty,1] \\[2pt]
% c(x-2)^2, & x\in[1,3] \\[2pt]
% d(x-2)^2+e(x-3)^3, & x\in[3,\infty)
% \end{cases}
% \]